% ============================================================
% Section: Layout Optimization (matches CAST_loss/voronoi_layout.py)
% 3-loss system: L_cap_res + L_cvt_norm + L_fea
% Core/Shell architecture with anisotropic power diagram
% ============================================================

\subsection{Hybrid Shape Initialization}
\label{subsec:shape_init}

Let $M$ denote the binary mask of the target container shape.
Standard centroidal Voronoi initialization~\cite{du1999centroidal} often fails to cover
thin protrusions in irregular silhouettes, producing degenerate
cells near extremities. To improve spatial coverage, we construct
an initialization set that combines skeleton-based anchors and
interior sampling.

\textbf{Skeleton Anchoring.}
We compute the Medial Axis Transform (MAT) of the mask $M$ and
extract endpoint vertices from the resulting skeleton graph. These
endpoints correspond to tips of narrow structural components
(ears, limbs, tails) and are selected as anchor sites $S_{\text{skel}}$,
subject to a minimum mutual distance filter to avoid clustering.
At most $\lfloor N/2 \rceil$ endpoints are retained so that the
remaining sites can be distributed volumetrically.

\textbf{Volumetric Filling.}
Additional sites $S_{\text{fill}}$ are sampled from local maxima of the
distance transform of $M$. These points densely populate interior
regions and stabilize subsequent layout optimization.

The final initialization set is
\[
S = S_{\text{skel}} \cup S_{\text{fill}},
\]
which provides a topology-aware starting configuration for layout
optimization under Voronoi-style spatial partitioning~\cite{du1999centroidal}.

% ============================================================
% Core/Shell Architecture
% ============================================================
\subsection{Core--Shell Spatial Priors}
\label{subsec:core_shell}

To prevent cell collapse in narrow or branching mask regions, we decompose the layout domain $\Omega$ into two complementary zones: a \emph{core} region with guaranteed per-cell ownership, and a \emph{shell} region where cells compete via the power diagram.

\textbf{Core construction.}
For each site $s_i$, we allocate a compact core region $C_i$ via connected multi-source region growing from $s_i$ within the foreground mask $M$. The core quota is set proportional to the average cell area:
\begin{equation}
q_i^{\text{core}} = \min\bigl(\alpha \cdot \bar{A},\; \beta \cdot c_i |\Omega|\bigr),
\end{equation}
where $\bar{A} = |\Omega|/N$ is the average cell area, $c_i$ is the target capacity of cell $i$, and $\alpha, \beta$ are scaling parameters. The cores $\{C_i\}$ are grown in parallel, ensuring disjoint ownership: each foreground pixel belongs to at most one core.

\textbf{Shell competition.}
The remaining foreground area $\Omega_{\text{free}} = \Omega \setminus \bigcup_i C_i$ constitutes the shell region. Each cell $i$ is assigned a local shell support mask $\mathcal{S}_i \subseteq \Omega_{\text{free}}$ based on proximity to its core. The soft power-diagram assignment operates \emph{only} within these shell support masks, while core pixels retain fixed ownership. This guarantees that every cell maintains a minimum viable area throughout optimization, preventing collapse even in narrow topology branches.

% ============================================================
% Optimization
% ============================================================
\subsection{Layout Optimization}
\label{subsec:layout_opt}

Given a set of $N$ selected keyframes, our goal is to arrange them within a layout domain $\Omega$ such that the resulting collage is both content-aware and geometrically coherent. The domain $\Omega$ corresponds to the foreground region of the layout mask with total area $|\Omega|$.

\vspace{1mm}
\noindent\textbf{Event-aware target allocation.}
For each keyframe $i$, we extract a primary event-support bounding box $B_i$ (Sec.~\ref{sec:key_event_abstraction}) with area $A_i^{\text{bbox}}$. We define a normalized target capacity:
\begin{equation}
c_i = \frac{\sqrt{A_i^{\text{bbox}}}}{\displaystyle\sum_{j=1}^{N} \sqrt{A_j^{\text{bbox}}}},
\end{equation}
where the square-root normalization prevents unusually large foreground boxes from dominating the layout, while retaining a monotonic relation between event extent and allocated area.

\vspace{1mm}
\noindent\textbf{Differentiable layout formulation.}
We model the layout as an anisotropic weighted power diagram parameterized by site positions $\{s_i\}_{i=1}^N$ and scalar weights $\{w_i\}_{i=1}^N$. Each cell $i$ is associated with an anisotropy matrix $A_i \in \mathbb{R}^{2 \times 2}$ derived from the target aspect ratio $\alpha_i^{\text{target}}$ of bounding box $B_i$. To avoid degenerate tessellation when all frames share the same aspect ratio, we normalize the matrices such that $\frac{1}{N}\sum_i A_i = I$. The anisotropic power distance from pixel $x$ to site $i$ is:
\begin{equation}
d_i(x) = (x - s_i)^\top A_i\,(x - s_i) - w_i.
\end{equation}

Soft cell assignment operates \emph{only on the shell region} $\Omega_{\text{free}}$. For each shell pixel $x \in \mathcal{S}_i$, we compute:
\begin{equation}
P_i^{\text{shell}}(x) = \frac{\exp\bigl(-\tau\, d_i(x)\bigr) \cdot \mathbb{1}_{\mathcal{S}_i}(x)}
              {\sum_{j=1}^{N} \exp\bigl(-\tau\, d_j(x)\bigr) \cdot \mathbb{1}_{\mathcal{S}_j}(x)},
\end{equation}
where $\tau$ is a temperature parameter and $\mathbb{1}_{\mathcal{S}_i}$ restricts competition to cells whose shell support covers pixel $x$. The full ownership probability combines core and shell:
\begin{equation}
P_i(x) = P_i^{\text{shell}}(x) + \mathbb{1}_{C_i}(x),
\end{equation}
where $\mathbb{1}_{C_i}(x) = 1$ if $x$ belongs to core $C_i$, ensuring guaranteed minimum ownership.

\vspace{1mm}
\noindent\textbf{Geometric attributes.}
The soft cell area decomposes into core and shell contributions:
\begin{equation}
A_i = A_i^{\text{core}} + A_i^{\text{shell}}, \quad
A_i^{\text{shell}} = \sum_{x \in \Omega_{\text{free}}} P_i^{\text{shell}}(x).
\end{equation}
The residual shell target for each cell is $T_i^{\text{res}} = \max(c_i |\Omega| - A_i^{\text{core}},\; 0)$. The effective area ratio is:
\begin{equation}
r_i = \frac{A_i^{\text{core}} + A_i^{\text{shell}}}{|\Omega|}.
\end{equation}
The soft centroid is computed over the full ownership:
\begin{equation}
\mu_i = \frac{\sum_{x \in \Omega} P_i(x)\,x}{A_i}.
\end{equation}

% ============================================================
% Loss Functions
% ============================================================
\vspace{1mm}
\noindent\textbf{Objective function.}
We optimize the layout by minimizing a composite loss:
\begin{equation}
\mathcal{L} =
\lambda_{\text{cap}}\, \gamma(t)\, \mathcal{L}_{\text{cap}} +
\lambda_{\text{cvt}}\, \mathcal{L}_{\text{cvt}} +
\lambda_{\text{fea}}\, \mathcal{L}_{\text{fea}},
\label{eq:total_loss}
\end{equation}
where $\gamma(t)$ is a smooth capacity ramp factor that gradually activates $\mathcal{L}_{\text{cap}}$ during a warmup phase (Sec.~\ref{subsec:schedule}). The three terms address residual shell capacity matching ($\mathcal{L}_{\text{cap}}$), centroidal regularity ($\mathcal{L}_{\text{cvt}}$), and feasibility enforcement ($\mathcal{L}_{\text{fea}}$), respectively.

An auxiliary collapse guard additionally penalizes cells whose shell-only ownership drops below a critical threshold, acting as a last-resort safety net independent of core support.

\begin{figure*}[t]
    \centering
    \begin{subfigure}[t]{0.3\linewidth}
        \centering
        \includegraphics[width=\linewidth]{images/loss_cap.png}
        \caption{Residual capacity: matches shell area to target after subtracting guaranteed core.}
    \end{subfigure}
    \hfill
    \begin{subfigure}[t]{0.3\linewidth}
        \centering
        \includegraphics[width=\linewidth]{images/loss_cvt.png}
        \caption{Centroidal regularity: pulls sites toward cell centroids, weighted by inverse cell size for small-cell focus.}
    \end{subfigure}
    \hfill
    \begin{subfigure}[t]{0.3\linewidth}
        \centering
        \includegraphics[width=\linewidth]{images/loss_fea.png}
        \caption{Feasibility: one-sided penalty that activates only when a cell's area falls below a safety floor.}
    \end{subfigure}
    \caption{
    Illustration of the proposed layout optimization losses.
    }
    \label{fig:loss_all}
    \vspace{-5pt}
\end{figure*}

\vspace{1mm}
\noindent\textbf{Residual Capacity Loss (Fig.~\ref{fig:loss_all}a).}
After subtracting core areas, we enforce that each cell's shell ownership matches its residual target. To ensure the loss magnitude is scale-invariant with respect to grid resolution, we normalize by the target itself:
\begin{equation}
\mathcal{L}_{\text{cap}} = \frac{1}{N}\sum_{i=1}^{N}
\left(\frac{A_i^{\text{shell}} - T_i^{\text{res}}}{T_i^{\text{res}} + \epsilon}\right)^{\!2},
\end{equation}
where $\epsilon$ is a small constant for numerical stability. This dimensionless formulation prevents the capacity term from dominating other losses, which would cause premature cell collapse in early iterations.

\vspace{1mm}
\noindent\textbf{Normalized Centroidal Loss (Fig.~\ref{fig:loss_all}b).}
To promote regular cell geometry, we encourage each site to coincide with its cell's centroid, with stronger penalization for smaller cells that are more vulnerable to collapse:
\begin{equation}
\mathcal{L}_{\text{cvt}} = \frac{1}{N}\sum_{i=1}^{N}
\frac{\|s_i - \mu_i\|^2}{\sqrt{\max(r_i,\; r_{\min})}},
\end{equation}
where $r_i$ is the effective area ratio and $r_{\min}$ is a clamping floor. The inverse square-root weighting focuses gradient signal on small or at-risk cells, stabilizing them without interfering with well-sized cells.

\vspace{1mm}
\noindent\textbf{Feasibility Penalty (Fig.~\ref{fig:loss_all}c).}
We define a one-sided penalty that activates only when a cell's effective area drops below a safety floor:
\begin{equation}
\mathcal{L}_{\text{fea}} = \phi(t) \sum_{i=1}^{N}
\left(\frac{[\,f_i - r_i\,]_+}{c_i + \epsilon}\right)^{\!2},
\end{equation}
where $f_i = 0.82\,c_i$ is the feasibility floor for cell $i$, $[\cdot]_+$ denotes the ReLU (zero for satisfied cells), $c_i$ normalizes the deficit so gradient magnitude is independent of cell count, and $\phi(t) = 1.5 - 0.5\,t$ is a time-decaying boost that applies stronger pressure in early iterations when cells are most vulnerable. Critically, this loss produces \emph{zero gradient} for cells at or above the floor, so large cells are never forced to shrink---which would cascade into collapse of adjacent cells in narrow mask regions.

% ============================================================
% Two-Stage Schedule
% ============================================================
\vspace{1mm}
\noindent\textbf{Two-stage optimization schedule.}
\label{subsec:schedule}
The optimization proceeds in two smoothly blended stages:

\emph{Stage~1 (Warmup):} During the initial $t < t_w$ fraction of iterations, the weight learning rate and capacity ramp factor $\gamma(t)$ start from small values and increase smoothly. This allows sites to settle into geometrically favorable positions under centroidal and feasibility guidance before shell capacity competition begins.

\emph{Stage~2 (Full optimization):} For $t \geq t_w$, all losses are fully active. The capacity ramp reaches $\gamma(t)=1$ and the weight learning rate reaches its maximum, enabling fine-grained area redistribution via additive power weights. The warmup ratio $t_w$ is adaptively extended when core support is large (high core fraction), since more guaranteed area requires gentler shell activation.

Additionally, a trust-region constraint limits site motion to a neighborhood of the initial anchor positions, preventing sites from drifting into distant mask components during optimization.

% ============================================================
% Rendering
% ============================================================
\subsection{Chronology-Preserving Rendering}
\label{subsec:composition}

\textbf{Chronological Assignment.}
To preserve temporal ordering, keyframes are assigned to layout
cells according to a strictly monotonic mapping $\pi(t)=j$,
following the broader motivation of story-driven visual summarization~\cite{lu2013story},
where $t$ denotes the temporal index of the keyframe and $j$ denotes
the spatial rank of the corresponding cell. Spatial ranks are determined
by a boustrophedon (zigzag row-major) ordering: cells are partitioned
into horizontal bands and sorted left-to-right in even rows and
right-to-left in odd rows, producing a continuous reading path through
the layout and reducing abrupt spatial jumps across neighboring rows.

\textbf{Saliency-Guided Warping.}
To fill highly irregular cell regions with minimal visible seams,
we deform each keyframe using a Thin-Plate Spline (TPS) warp~\cite{bookstein2002principal} that
preserves salient foreground geometry. A saliency map derived from
the event-support mask assigns high rigidity weights to foreground
control points and low weights to background control points,
concentrating deformation in non-salient regions while keeping
the primary subject intact.

% ============================================================
% CITE LIST
% ============================================================
%
% \cite{du1999centroidal}
%   Du, Q., Faber, V., Gunzburger, M.
%   "Centroidal Voronoi tessellations: Applications and algorithms."
%   SIAM Review, 41(4):637-676, 1999.
%
% \cite{bookstein2002principal}
%   Bookstein, F. L.
%   "Principal warps: Thin-plate splines and the decomposition
%    of deformations."
%   IEEE TPAMI, 11(6):567-585, 1989.
%
% \cite{lu2013story}
%   Lu, Z., Grauman, K.
%   "Story-driven summarization for egocentric video."
%   IEEE CVPR, 2013.